\section{The countable construction and the actual center}
\label{sec:countable-construction}

We now apply Theorem~\ref{thm:lifting} repeatedly.  Recall its
full conclusion: from a finite-dimensional central division algebra $A/k$
and a tuple of free series $X$ whose constant term generates $A$ and whose
two linear coefficients are independent modulo inner directions, it
produces such a tuple over a countable extension $k'$ and a central
division algebra of strictly larger degree.  It preserves every old
polynomial identity and gives a monic algebraic relation for any chosen
series $w$.  The countability, generation, and two-direction conclusions
are all needed to iterate this theorem.

The resulting division algebras need not embed in one another.  We use
an ultraproduct to retain their polynomial relations, and then prove
separately that the resulting algebra has infinite dimension over its
actual center.

\subsection{Scheduling all polynomial expressions}

Start with
\[
 k_0=\overline{\mathbb Q},\qquad A_0=k_0,\qquad
 X^{(0)}=(z_1,z_2).
\]
Its constant pair $(0,0)$ generates $k_0$ as a unital $k_0$-algebra.
The inner-direction space is zero, and its two linear coefficients
$(1,0)$ and $(0,1)$ are independent.  Thus it satisfies the hypotheses
of Theorem~\ref{thm:lifting}.

Inductively we obtain countable fields
\[
 k_0\subseteq k_1\subseteq k_2\subseteq\cdots,
\]
finite-dimensional central division algebras $A_i/k_i$ of strictly
increasing degrees, and pairs
\[
 X^{(i)}\in A_i\langle\!\langle z_1,z_2\rangle\!\rangle^2
\]
satisfying the same generation and transverse-independence hypotheses.
Write $(x_i,y_i)=X^{(i)}_0$ for the constant pair; in particular,
\begin{equation}
 A_i=k_i[x_i,y_i].
 \label{eq:stage-generation}
\end{equation}

Here is a schedule that ensures algebraicity of every expression
needed in the limit.  When $k_i$ has been constructed, enumerate the
countable set $k_i\langle x,y\rangle$ as
$(p_{i,j})_{j\geq0}$.  Enumerate the pairs $(i,j)\in\mathbb N^2$ by
increasing $i+j$, in a fixed order within each diagonal.  If $(i_n,j_n)$
is the $n$th pair, beginning at $n=0$, then $i_n\leq n$.
At stage $n$, apply the lifting theorem to
\[
 w=p_{i_n,j_n}(X^{(n)}).
\]
This expression is defined over $k_n$, since $k_{i_n}\subseteq k_n$.
The identity $w-p_{i_n,j_n}(X^{(n)})=0$ is preserved, so the new relation
for the lifted $w$ gives a nonzero monic polynomial
\[
 q_{i_n,j_n}(t)\in k_{n+1}[t]
\]
such that
\begin{equation}
 q_{i_n,j_n}\bigl(p_{i_n,j_n}(X^{(n+1)})\bigr)=0.
 \label{eq:retained-formal-relation}
\end{equation}
Every later lifting preserves this polynomial identity, now regarded
over its coefficient field $k_{n+1}$.  Taking constant terms shows that
the corresponding identity holds at $(x_h,y_h)$ for every $h\geq n+1$.
Each polynomial over each $k_i$ is treated at a finite stage.

\subsection{A two-generated algebraic division ring}

Fix a nonprincipal ultrafilter $\mathcal U$ on $\mathbb N$ and form
(see \cite{Keisler2010} for the general construction)
\[
 \mathcal D=\prod_{i\geq0} A_i\big/\mathcal U.
\]
Here two sequences represent the same element when their equality set
belongs to $\mathcal U$.  This is a division ring: a nonzero class has
nonzero coordinates on a set in $\mathcal U$, and taking their
coordinatewise inverses on that set gives its inverse.

Put $k_\infty=\bigcup_i k_i$.  An element $c\in k_j$ determines a class
in $\mathcal D$ by taking its image in the central field $k_i\subset A_i$
for $i\geq j$ and choosing arbitrary coordinates at the finitely many
earlier indices.  These classes are independent of the early choices
and of $j$.  They define a central field embedding
\[
 k_\infty\hookrightarrow\mathcal D:
\]
the field operations agree at every sufficiently large index, and a
nonzero field element stays nonzero at all those indices.  This
construction uses only the inclusions of coefficient fields; it does
not require embeddings between the algebras $A_i$.

Let
\[
 x=[(x_i)_i],\qquad y=[(y_i)_i],\qquad
 R=k_\infty[x,y]\subseteq\mathcal D.
\]
Every $r\in R$ has the form $p(x,y)$ for some
$p\in k_j\langle x,y\rangle$.  The schedule and
\eqref{eq:retained-formal-relation} provide a nonzero monic polynomial
$q\in k_h[t]$, for a suitable $h$, such that
$q(p(x_i,y_i))=0$ at every sufficiently large index.  Therefore
$q(r)=0$ in $\mathcal D$.  Every element of $R$ is algebraic over
$k_\infty$.

This already makes $R$ a division ring.  For $0\ne r\in R$, choose a
nonzero relation of least degree
\[
 a_0+a_1r+\cdots+a_nr^n=0,
 \qquad a_j\in k_\infty.
\]
Its constant coefficient is nonzero: otherwise multiplication by
$r^{-1}$ in the ambient division ring $\mathcal D$ would give a shorter
relation.  Consequently
\[
 r^{-1}=-a_0^{-1}(a_1+a_2r+\cdots+a_nr^{n-1})\in R.
\]
Since $k_\infty$ is countable and $R$ is generated by two elements as an
algebra, $R$ is countable.  Its characteristic is zero.

\subsection{Infinite dimension and passage to the center}

We first prove
\begin{equation}
 [R:k_\infty]=\infty.
 \label{eq:infinite-auxiliary-dimension}
\end{equation}
For each $i$ and $\ell\geq0$, let $V_{i,\ell}$ be the $k_i$-span of the
words in $x_i,y_i$ of length at most $\ell$, including the empty word
$1$.  If $V_{i,\ell}=V_{i,\ell+1}$, then $V_{i,\ell}$ is stable under
right multiplication by $x_i$ and $y_i$ and contains every word.
By \eqref{eq:stage-generation}, it equals $A_i$.  Until this happens,
the dimensions increase strictly, starting from $1$.

Fix $\ell$.  The dimensions $[A_i:k_i]$ tend to infinity, so
\[
 \dim_{k_i}V_{i,\ell}\geq\ell+1
\]
for every sufficiently large $i$.  There are only finitely many
subsets of size $\ell+1$ in the finite list of words of length at most
$\ell$.  The ultrafilter therefore selects one fixed subset
$w_0,\ldots,w_\ell$ that is linearly independent at a set of stages
belonging to $\mathcal U$.

These same words are independent over the embedded $k_\infty$ in $R$.
Indeed, the coefficients of a putative relation
\[
 \sum_{j=0}^{\ell} c_j w_j(x,y)=0
\]
all belong to one field $k_h$.  Equality in the ultraproduct holds on
a set in $\mathcal U$.  Intersecting it with the independence set and
the tail $i\geq h$ gives a nonempty set where stage independence
forces every $c_j$ to vanish.  Thus $R$ has arbitrarily large finite
linearly independent sets, proving
\eqref{eq:infinite-auxiliary-dimension}.

Finally let $F=Z(R)$ be the actual center.  It contains $k_\infty$ and
is algebraic over $k_\infty$, since every element of $R$ is.  In
particular, every element of $R$ is algebraic over $F$.  Also
\[
 R=k_\infty[x,y]=F[x,y]=F(x,y),
\]
because $F\subseteq R$ and $R$ is already a division ring.

We use finite descent of the multiplication table and generator
coordinates; compare the proof of the local-finiteness criterion in
\cite[Section~4]{DeoBienHai2019}. Suppose, for a contradiction, that
$[R:F]=n<\infty$. Choose an
$F$-basis $b_1=1,b_2,\ldots,b_n$ of $R$, and write
\[
 b_i b_j=\sum_{h=1}^n c_{ijh}b_h,
 \qquad
 x=\sum_{h=1}^n u_hb_h,
 \qquad
 y=\sum_{h=1}^n v_hb_h,
\]
with all coefficients in $F$.  Let $E$ be the subfield of $F$
generated over $k_\infty$ by these finitely many coefficients.
Algebraicity gives $[E:k_\infty]<\infty$.  The space
\[
 B=\sum_{h=1}^n E b_h
\]
contains $1,x,y$ and is closed under multiplication by the displayed
structure constants.  It is therefore a $k_\infty$-subalgebra of $R$
containing $k_\infty[x,y]=R$, so $B=R$.  This implies
\[
 [R:k_\infty]\leq n[E:k_\infty]<\infty,
\]
contrary to \eqref{eq:infinite-auxiliary-dimension}.

Thus the countable characteristic-zero division ring $D=R$ is
algebraic over its actual center $F$, while
\[
 D=F[x,y]=F(x,y),\qquad [D:F]=\infty.
\]
The finite set $\{x,y\}$ witnesses the failure of local finiteness.

This proves Theorem~\ref{thm:main}, subject only to Theorem~\ref{thm:lifting},
which we now prove.
